\section{Data construction details}\label{app:data}

\paragraph{Repricing distribution.} For each year-end $t$ and each liability $i$ with par (or balance) $w_i$ and years to repricing $\tau_i$, the repricing distribution is $F_t(h)=\sum_i w_i\,\mathbf 1\{\tau_i\le h\}/\sum_i w_i$. Repricing dates are maturity for bills, notes, bonds and TIPS; one week for floating-rate notes; and zero (overnight) for reserve balances and reverse repurchase agreements. The clock adds growth-financing issuance at the new rate, $P_t(h)=1-(1-F_t(h))e^{-gh}$, with $g=4\%$ for the descriptive clock (Figure~\ref{fig:clock}) and the year's $g$ in the limit calculations. The erosion clock $P_N$ is built the same way from non-indexed liabilities only.

\paragraph{Zero-interest base.} Currency in circulation is from H.4.1 (series RESTBC), in the same week as reserves, from 2003. Before that we use December averages of currency in circulation and reserve balances (FRED CURRCIR and RESBALNS); the two currency sources differ by about 1\% where they overlap. Reserves are in the zero-interest base through 2007 and in overnight interest-bearing debt from 2008.

\paragraph{Treasury-only and consolidated stocks.} The Treasury-only stock is all unmatured marketable Treasuries in MSPD Table 3 at the year-end, one row per CUSIP, including those held by the Federal Reserve. The consolidated stock is marketable Treasuries minus SOMA holdings, matched by CUSIP (with TIPS inflation compensation), plus reserve balances and reverse repos from the H.4.1 week closest to, and not after, the year-end. SOMA CUSIPs that are absent from the MSPD year-end snapshot all mature between the SOMA date and 31 December (at most \$61bn, 0.3\% of the stock) and are dropped.

\paragraph{1980--2002.} FD-5 reports privately held marketable debt in five remaining-maturity buckets: within 1 year, 1--5, 5--10, 10--20, and 20 years and over. Within each bucket we spread par uniformly (upper edge of the last bucket: 30 years). On the exact security-level data for 2003--07, this bucketing understates the ten-year inflation-layer ratio by 0.077 (s.d.\ 0.016), $\int P$ by 0.167, and the remaining-maturity rate $r$ by 0.28pp; we add these corrections to the 1980--2002 values. TIPS (1997--2002) are removed from the erosion base by bucket: their December totals come from Treasury Bulletin table FD-2 and their maturity profile from the TIPS outstanding at each year-end. Before 2003 this also removes the Fed's small TIPS holdings (8\% of TIPS in 2003); scaling the removal by the private share instead changes the ratio by at most 0.006.

\paragraph{Required inflation.} The integrals in Propositions~\ref{prop:price} and~\ref{prop:kappa} are evaluated by the trapezoid rule on a monthly grid over $[0,15]$ years, with $F(h)$ interpolated between the sorted repricing dates of the security-level data.

\paragraph{Backtest projection (Section~\ref{sec:backtest}).} From the year-end portfolio, every surviving security keeps its rate: the coupon for notes and bonds, the real coupon for TIPS, and the three-month rate plus spread for floating-rate notes. Each month, maturing principal plus net new borrowing is reissued in year $t$'s gross auction mix, at that month's H.15 yield for the tenor. The projected average rate is the par-weighted rate of the resulting portfolio, and we score its change against the change in FiscalData's official average interest rate on total marketable debt. ``Realized borrowing'' interpolates the actual path of marketable debt; ``trend borrowing'' grows the year-$t$ stock at 4\% a year.

\paragraph{Fiscal-reaction coefficient.} \citet{auerbach2024} estimate the legislated primary-surplus response to lagged net interest. Their coefficient is in dollars of surplus per dollar of interest, which is the marginal offset $\varphi$ of Section~\ref{sec:model}. \citet{eichengreen2026} instead estimate a log-log elasticity of the surplus to debt service, which maps to $\varphi$ only through $s/(\bar r b)$; that ratio is undefined when the primary balance is negative, so we use it as qualitative evidence only.

\paragraph{United Kingdom (Section~\ref{sec:international}).} We use the DMO's report D1A for the last business day of each year, 2007--2025. Conventional gilts, including ``rump'' gilts, enter at their nominal amount; index-linked gilts, with both the three-month and the eight-month indexation lag, enter at their nominal amount including the inflation uplift. Undated gilts have no repricing date and are given one of 100 years until the Treasury announced their redemption in 2014--15. The parsed sums reproduce the total stated in each report. APF holdings of each gilt at any date are the sum of the nominal amounts in all operations settled by that date: APF purchases in 2009--21, including reinvestments; the 2022 financial-stability purchases of long-dated conventional and index-linked gilts; minus active sales from November 2022 and the financial-stability sales of November 2022 to January 2023. A holding drops out at the gilt's redemption, read from the Bank's bond code, and index-linked holdings are uplifted with the DMO ratio of amount to nominal for the same gilt. Reserve balances and notes and coin are the last December observations of Bank of England series LPMBL22 and LPMAVAA, and GDP is the calendar-year sum of quarterly nominal GDP (FRED UKNGDP). The marginal cost $r$ prices each gilt at the calendar-year average yield for its original tenor, interpolated between SONIA and the Bank's 5-, 10- and 20-year nominal par yields and flat beyond 20 years; undated gilts are priced at 20 years. Real growth is from FRED NGDPRSAXDCGBQ. The inflation test uses the Bank's 5-, 10- and 20-year zero-coupon nominal, real and implied-inflation curves (monthly averages), with Bank Rate at zero maturity and the real curve flat below five years, and ONS series D7BT (CPI) and CHAW (RPI).

\paragraph{Japan (Section~\ref{sec:international}).} JGBs by issue at each fiscal year-end come from Table 34 of the Ministry of Finance's debt yearbook (general bonds and FILP bonds). Columns are located by their header labels, because the layout differs across editions, and issues are classified as index-linked, floating-rate (repricing at the next six-month reset), discount bills and fixed-rate. BoJ holdings are matched by issue name and number from the BoJ's release of JGB holdings by issue at the fiscal year-end (the last business day of March). Financing bills are the treasury-discount-bill total in the BoJ's public-finance statistics (PF02) minus the treasury bills in Table 34, and reprice at three months; BoJ holdings of bills (balance-sheet item MABJMA5A) and government holdings of bills and JGBs (PF02) are allocated pro rata. Current accounts are the end-March balance (MABJML11). Through February 2024, the 0\% tier's share of current accounts is the March average share of balances at a zero rate among the three tiers (MD08), applied to the end-March balance; after the March 2024 reform the zero-rate part is required reserves (MD07, average outstanding). Banknotes (MABJML1) complete the zero-interest base. GDP is the fiscal-year average of quarterly nominal GDP at annual rates (FRED JPNNGDP). The marginal cost $r$ prices each issue at the fiscal-year average of the Ministry of Finance's constant-maturity yield for its original tenor; bills and floating-rate bonds are priced at one year and index-linked bonds at the ten-year nominal yield. Real growth is the fiscal-year growth of FRED JPNRGDPEXP. The inflation test uses monthly averages of the constant-maturity JGB yields, with the call rate (BoJ series STRDCLUCON) at zero maturity, and the current-account tier balances of MD08 to weight the rates of $+0.1\%$, $0\%$ and $-0.1\%$. CPI (all items) is the Statistics Bureau's 2020-base series from January 2020, chained to the OECD series (FRED JPNCPIALLMINMEI) before that.
